\documentclass[a4paper,12pt,fleqn]{article}%fleqn pour aligner les equations à gauche
\usepackage{/home/rweye/Documents/Overleaf/Styles/Cours}

\title{\vspace{-1cm}\large 
	\begin{tabular}{|>{\centering}m{5cm}|>{\centering}m{8cm}|>{\centering\arraybackslash}m{4cm}|}
		\hline
		Activité 2 & \textbf{\textsc{La, mais à distance !}} &  Chapitre 6 \\
		\hline
\end{tabular}
\vspace{-5em}}
\date{}

\begin{document}
\begin{tabular}{|>{\centering}m{5cm}|>{\centering}m{8cm}|>{\centering\arraybackslash}m{4cm}|}
		\hline
		Activité 2 & \textbf{\textsc{La, mais à distance !}} &  Chapitre 6 \\
		\hline
\end{tabular}\\

%\vspace{1em}
Lors de la période de confinement les habitudes de travail ont du s'adapter aux conditions du distanciel. Au printemps dernier, un pianiste devait faire accorder don instrument comme il le fait régulièrement. Pendant la période de confinement, l’artisan qu’il a l’habitude de consulter ne pouvait pas faire le déplacement. Mais ne voulant pas mettre son client dans l’embarras il lui propose de procéder lui-même à son accordage et lui demande, pour évaluer si l’instrument est désaccordé, de lui envoyer un enregistrement d’une touche (le La3). Le pianiste lui envoie donc un fichier “La3.wav” que l’artisan peut alors comparer au La3 produit par son diapason, fichier “diapason.wav”.\\
\textbf{En vous mettant à la place de l’accordeur de piano. Vérifier que la touche du piano est bien accordée.}\\

\vspace{1em}
\begin{tabular}{|>{\centering}m{5cm}|>{\centering}m{8cm}|>{\centering\arraybackslash}m{4cm}|}
		\hline
		Activité 2 & \textbf{\textsc{La, mais à distance !}} &  Chapitre 6 \\
		\hline
\end{tabular}\\

%\vspace{1em}
Lors de la période de confinement les habitudes de travail ont du s'adapter aux conditions du distanciel. Au printemps dernier, un pianiste devait faire accorder don instrument comme il le fait régulièrement. Pendant la période de confinement, l’artisan qu’il a l’habitude de consulter ne pouvait pas faire le déplacement. Mais ne voulant pas mettre son client dans l’embarras il lui propose de procéder lui-même à son accordage et lui demande, pour évaluer si l’instrument est désaccordé, de lui envoyer un enregistrement d’une touche (le La3). Le pianiste lui envoie donc un fichier “La3.wav” que l’artisan peut alors comparer au La3 produit par son diapason, fichier “diapason.wav”.\\
\textbf{En vous mettant à la place de l’accordeur de piano. Vérifier que la touche du piano est bien accordée.}\\
\vspace{1em}

\begin{tabular}{|>{\centering}m{5cm}|>{\centering}m{8cm}|>{\centering\arraybackslash}m{4cm}|}
		\hline
		Activité 2 & \textbf{\textsc{La, mais à distance !}} &  Chapitre 6 \\
		\hline
\end{tabular}\\

%\vspace{1em}
Lors de la période de confinement les habitudes de travail ont du s'adapter aux conditions du distanciel. Au printemps dernier, un pianiste devait faire accorder don instrument comme il le fait régulièrement. Pendant la période de confinement, l’artisan qu’il a l’habitude de consulter ne pouvait pas faire le déplacement. Mais ne voulant pas mettre son client dans l’embarras il lui propose de procéder lui-même à son accordage et lui demande, pour évaluer si l’instrument est désaccordé, de lui envoyer un enregistrement d’une touche (le La3). Le pianiste lui envoie donc un fichier “La3.wav” que l’artisan peut alors comparer au La3 produit par son diapason, fichier “diapason.wav”.\\
\textbf{En vous mettant à la place de l’accordeur de piano. Vérifier que la touche du piano est bien accordée.}\\
\vspace{1em}

\begin{tabular}{|>{\centering}m{5cm}|>{\centering}m{8cm}|>{\centering\arraybackslash}m{4cm}|}
		\hline
		Activité 2 & \textbf{\textsc{La, mais à distance !}} &  Chapitre 6 \\
		\hline
\end{tabular}\\

%
Lors de la période de confinement les habitudes de travail ont du s'adapter aux conditions du distanciel. Au printemps dernier, un pianiste devait faire accorder don instrument comme il le fait régulièrement. Pendant la période de confinement, l’artisan qu’il a l’habitude de consulter ne pouvait pas faire le déplacement. Mais ne voulant pas mettre son client dans l’embarras il lui propose de procéder lui-même à son accordage et lui demande, pour évaluer si l’instrument est désaccordé, de lui envoyer un enregistrement d’une touche (le La3). Le pianiste lui envoie donc un fichier “La3.wav” que l’artisan peut alors comparer au La3 produit par son diapason, fichier “diapason.wav”.\\
\textbf{En vous mettant à la place de l’accordeur de piano. Vérifier que la touche du piano est bien accordée.}\\


\end{document}